% The Master Equation of SIGIL — companion note v0.2
% LaTeX rendition of SIGIL_MASTER_EQUATION_v0.md (v0 2026-05-30, v0.1 addendum
% 2026-05-31), corrected per the 2026-07-18 external review:
%  - framed throughout as a CONJECTURAL EFFECTIVE ACTION (correspondences, not
%    derivations)
%  - PID stability discriminant corrected (K_p^2 >= 4 K_d K_i)
%  - gauge-healing language fixed (re-synchronization dynamics, not a gauge
%    transformation)
%  - knot formulation fixed (braid closure -> link; invariants are one-way)
%  - crypto composition stated per-property
% Build: pdflatex x2.
\documentclass[10pt,a4paper]{article}

\usepackage[margin=2.6cm]{geometry}
\usepackage[T1]{fontenc}
\usepackage{lmodern}
\usepackage{amsmath,amssymb,amsfonts}
\usepackage{xcolor}
\usepackage{booktabs}
\usepackage{enumitem}
\usepackage{microtype}
\usepackage[colorlinks=true,linkcolor=sigilviolet,urlcolor=sigilcyan,citecolor=sigilviolet]{hyperref}
\usepackage{xurl}

\definecolor{sigilviolet}{HTML}{8B5CF6}
\definecolor{sigilgold}{HTML}{B7791F}
\definecolor{sigilcyan}{HTML}{0E7490}
\definecolor{ink}{HTML}{1A1428}

\hypersetup{
  pdftitle={The Master Equation of SIGIL — a Conjectural Effective Action for a BlockDAG (v0.2)},
  pdfauthor={bitknight.dipper688@passmail.net},
  pdfsubject={Companion note to the SIGIL whitepaper v1.1; sigilgraph repo, branch feat/commons-tithe-routing-0.36.1},
  pdfkeywords={SIGIL, BlockDAG, effective action, variational principle, knot theory, post-quantum}
}

\newcommand{\SSIGIL}{\mathcal{S}_{\mathrm{SIGIL}}}
\newcommand{\pretend}[1]{\textcolor{sigilgold}{\textit{#1}}}

\title{\textcolor{ink}{\LARGE The Master Equation of SIGIL}\\[5pt]
\large A Conjectural Effective Action for a BlockDAG\\[4pt]
\normalsize Companion note v0.2 · 2026-07-18 · supersedes v0/v0.1 (2026-05-30/31)}
\author{\texttt{bitknight.dipper688@passmail.net}\\[2pt]
\normalsize companion to \emph{SIGIL — The Variational Chain}, whitepaper v1.1}
\date{}

\begin{document}
\maketitle

\begin{abstract}
\noindent
This note states the conjectural effective action that organizes the SIGIL protocol,
defines its fields, maps each variation to a shipped or staged protocol component, and
lists the falsifiable predictions the correspondence generates. \textbf{Its epistemic
status is stated up front: this is a design principle and research programme, not
accomplished mathematics.} The component Lagrangians are unspecified beyond their field
content; no discretization of the BlockDAG into a lattice field theory has been
constructed; each ``field equation'' below is a correspondence between a standard
equation and a mechanism, adopted because it forces holistic design and yields testable
predictions. One prediction generated under this discipline — the gossip delivery law —
has completed the full arc from hypothesis to measured closed form. v0.2 incorporates
corrections from external review: the PID stability discriminant, the gauge-healing
language, the knot formulation, and the per-property statement of cryptographic
composition.
\end{abstract}

\section{The equation}

SIGIL's BlockDAG braid $\mathcal{G}$ — vertices are blocks, edges are causal links — is
modeled as a Lorentzian 4-manifold. We conjecture its dynamics are organized by one
variational principle:

\begin{equation}
\boxed{\;\delta \SSIGIL[g,\,A,\,\varphi,\,J,\,K,\,\Sigma] \;=\; 0\;}
\end{equation}

\begin{equation}
\SSIGIL = \int_{\mathcal{G}} d^{4}\sigma\, \sqrt{|g|}\,
\Bigl[
\tfrac{1}{2\kappa}\bigl(R - 2\Lambda_{\mathrm{emission}}(\varphi)\bigr)
+ \mathcal{L}_{\mathrm{settle}}(J,g)
+ \mathcal{L}_{\mathrm{witness}}(A,g)
+ \mathcal{L}_{\mathrm{topo}}(K)
+ \mathcal{L}_{\mathrm{crypto}}(\Sigma,g)
\Bigr]
\end{equation}

The claim this note is entitled to make: every SIGIL mechanism must identify which term
it belongs to, and a proposed mechanism touching two or more terms is a design smell.
The claim it is \emph{not} yet entitled to make: that any mechanism has been derived
from the action. The $\mathcal{L}$'s are names with declared field content, awaiting
explicit form.

\section{The fields}

\begin{center}
\small
\begin{tabular}{@{}llp{8.8cm}@{}}
\toprule
Symbol & Type & Protocol meaning \\
\midrule
$\mathcal{G}$ & Lorentzian 4-manifold & The discrete BlockDAG braid, conjecturally
modeled as a Lorentzian 4-manifold; the fabric of the chain. \\
$g_{\mu\nu}$ & metric tensor & Conjectured to be induced by agent-reputation density;
reputation as curvature. \pretend{Heuristic — functional form unfit.} \\
$R$ & Ricci scalar & Consensus tension; forks raise it, convergence minimizes it. \\
$\Lambda_{\mathrm{emission}}(\varphi)$ & scalar field & The emission ``cosmological
constant,'' PID-controlled toward target supply. \\
$A_\mu$ & gauge connection & Witness attestations; $F = dA + A\wedge A$ is the
witness-disagreement field. \\
$J^\mu$ & conserved current & Settlement flow; $J^0$ value density, $J^i$ settlement
current. \\
$K$ & link (closed braid) & The link obtained by closing the DAG's braid over a bounded
window; consensus-equivalent histories share its invariants. \\
$\Sigma = (\Sigma_{\mathrm{iso}}, \Sigma_{\mathrm{lat}}, \Sigma_{\mathrm{hash}})$ &
security connection & Three disjoint post-quantum hardness assumptions, one per layer
they guard. \\
$\kappa = 8\pi G_{\mathrm{agent}}$ & coupling & The gravitational constant of the agent
economy; empirical. \\
\bottomrule
\end{tabular}
\end{center}

\section{The correspondences}

Each variation names one protocol component. ``Names,'' not ``yields'': see the abstract.

\subsection{$\delta\SSIGIL/\delta J^\mu = 0$ — conservation (the ledger)}

\begin{equation}
\nabla_\mu J^\mu = 0 \;\Longleftrightarrow\;
\partial_t \rho_{\mathrm{value}} + \nabla\!\cdot\!\mathbf{J}_{\mathrm{settle}}
= S_{\mathrm{emission}}(\varphi)
\end{equation}

Value conservation with the emission source term. Implemented as the
\texttt{commit\_state\_transition} chokepoint — the only code path permitted to mutate
chain state — with per-delta conservation checks and the 21M cap as the integral form.

\subsection{$\delta\SSIGIL/\delta\varphi = 0$ — Klein--Gordon emission (monetary policy)}

\begin{equation}
\ddot{\varphi} + 3H\dot{\varphi} + V'(\varphi) = 0,\qquad
V(\varphi) = K_p\,e^2 + K_i\,e\!\!\int\! e\,dt + K_d\,(\dot e)^2
\end{equation}

The PID emission controller cast as a scalar field in a potential; Hubble friction is
network expansion. The displayed $V$ is \emph{schematic}: $e$, $\int e\,dt$, and
$\dot e$ are controller states, not functions of $\varphi$ alone, so $V$ is not an
ordinary local potential — a faithful treatment introduces auxiliary fields (or defines
$e(\varphi)$ explicitly). The mapping is exact only for quadratic potentials —
anti-windup,
derivative filtering, and gain scheduling break the analogy and are deferred to
higher-order correction terms. \textbf{Corrected stability condition (v0.2):} from the
linearized characteristic equation $K_d s^2 + K_p s + K_i = 0$, the non-oscillatory
(overdamped) condition is
\begin{equation}
K_p^2 \;\ge\; 4\,K_d K_i
\end{equation}
v0.1 stated $K_d^2 > 4K_pK_i$, interchanging the damping and inertia roles; a fresh
derivation against the \emph{implemented} controller (which is not exactly this
second-order plant) is part of the prediction's test protocol.

\subsection{$\delta\SSIGIL/\delta K = 0$ — topological invariance (fork detection)}

\begin{equation}
\frac{d}{dt}\,V_K(q)\Big|_{\mathrm{consensus}} = 0
\end{equation}

where $K$ is the link obtained by closing the braid over a bounded window (v0.1's
$K \in \pi_1(\mathcal{G})$ was a category error — a knot is not an element of the
fundamental group). Valid consensus moves are Reidemeister moves; invariants are
conserved; a fork \emph{may} change the bounded-window invariant — when it does,
divergence is proven, and equality of invariants never proves equality of histories. Two
constraints bind the programme. \emph{Computational:} the Jones polynomial is \#P-hard in general, so the
per-block commitment uses cheap invariants — linking number ($O(\mathrm{crossings})$) or
the Alexander polynomial (polynomial-time via the Burau determinant) — over a bounded
window, with Jones/Khovanov machinery reserved for rare deep-fork adjudication.
\emph{Logical:} cheap invariants fingerprint, they do not classify. A differing invariant
proves divergence; identical invariants do not prove identical histories. The topology
commitment is a one-way tripwire that complements, never replaces, the state-root
comparison.

\subsection{$\delta\SSIGIL/\delta A^\mu = 0$ — Yang--Mills witnesses (divergence
detection)}

\begin{equation}
\nabla_\nu F^{\mu\nu} = J^\mu_{\mathrm{witness}},\qquad F = dA + A\wedge A
\end{equation}

Diverging state-root tracks make the curvature $F$ locally non-zero; detection is the
measurement of $|F|$ above threshold. \textbf{Corrected language (v0.2):} healing is
\emph{re-synchronization dynamics} that drives $F$ back toward zero. It is not a gauge
transformation — curvature is gauge-covariant, so a gauge transformation relabels the
connection without repairing the underlying disagreement. In code: four state roots per
header, independent recomputation at every node, self-halt on mismatch.

\subsection{$\delta\SSIGIL/\delta\Sigma = 0$ — crypto-agility (the security floor)}

Three disjoint post-quantum hardness families, each guarding the layer where its
strengths bind: $\Sigma_{\mathrm{iso}}$ (isogeny — SQIsign L5; identity and provenance,
sign once, 292\,B), $\Sigma_{\mathrm{lat}}$ (lattice — ML-DSA/ML-KEM; per-transaction
and transport, fast batched verification), $\Sigma_{\mathrm{hash}}$ (FRI-STARK and the
Ajtai/SIS fold; state and tip-proof, succinct whole-chain verification, no trusted
setup). The correspondence motivates the dispatch rule: retire a leg the instant its
hardness margin drops; the others carry what they guard. \textbf{Composition, stated carefully
(v0.2):} because the legs protect \emph{different} properties, the security of any single
property is the floor of the primitives guarding \emph{that} property; an adversary
targeting one property attacks one leg. The break-all-three bar applies only to forging a
block's complete commitment stack as a coherent whole, and a formal joint security model
is open work.

\subsection{The constraint surface — the wall}

SIGIL imposes the engineering feasibility condition $\Gamma_{\mathrm{verify}} \ge
\Gamma_{\mathrm{state}}$; we call the boundary ``the wall.'' It is an operational
constraint, not a proven condition for the existence of an action extremum. Measured
2026-05-30 at $\sim$1840:1 against verification, resolved not by faster per-signature
verification but by changing
$\Sigma_{\mathrm{hash}}$: the fold collapsed whole-chain verification to a constant-size
check (2{,}568\,B; 342\,ms full verification at the 100k-block benchmark; 3\,$\mu$s
live incremental tip check). The wall was a cryptographic choice, not a hardware limit —
the one structural claim of this note that events have since borne out.

\section{Predictions}

The correspondence earns its keep only by predicting. Status at v0.2:

\begin{center}
\small
\begin{tabular}{@{}p{7.2cm}lp{5.2cm}@{}}
\toprule
Prediction & Status & Test \\
\midrule
1. Fork survival $\propto \exp(-\Delta R\,\tau_{\mathrm{block}})$ & \pretend{open} &
fork-survival statistics vs.\ producer reputation \\
2. PID non-oscillatory iff $K_p^2 \ge 4K_dK_i$ (corrected) & \pretend{open} & derive
against implemented controller; tune to critical value \\
3. Bounded-window topology commitment detects order-tampering forks &
partially met & braid gates: 16/16 permutations agree, 5/5 tamper rejected; header
commitment staged \\
4. The verification wall scales with $\Sigma$, not hardware & confirmed in kind & the
fold moved the wall; SIMD did not \\
5. Bulk/boundary duality (settlement throughput $\sim$ boundary entropy production) &
\pretend{open} & compare vs.\ reputation-graph entropy \\
6. Validator-departure entropy law (Hawking analogue) & \pretend{open} & tip-proof rate
vs.\ validator-set entropy \\
7. Gossip delivery follows $D(p,r)=1-p^r$ & \textbf{confirmed} & 22 seeded simulations,
RMS 0.4\,pp; shipped as the $r{=}3$ provisioning default \\
\bottomrule
\end{tabular}
\end{center}

Prediction 7 is the template: hypothesis $\to$ seeded simulation campaign $\to$ closed
form $\to$ residual $\to$ operator-facing rule $\to$ published paper. Predictions 1--6
are held to that arc.

\section{What's still pretend}

In one place, so no reader has to hunt: (i) the action is conjectural — Lagrangians
unspecified, no lattice discretization, correspondences not derivations; (ii) the
reputation metric is a heuristic awaiting a fit to measured consensus drift; (iii) the
tractable topology programme is a fingerprint, not a classifier, and the header
commitment is staged; (iv) AdS/CFT on a discrete graph is analogy pending rigorous
discretization; (v) the PID/Klein--Gordon mapping holds only for quadratic potentials;
(vi) per-property crypto floors are design intent — the joint composition bound and the
fold's concrete Ajtai/SIS soundness reduction are open; (vii) the 1840:1 saddle
measurement was taken on one 48-core host — the saddle condition is hardware-agnostic,
its scale is not.

\section{Method}

This note was produced under the project's measure-first discipline: the fields were
chosen to match modules that already exist in the workspace (\texttt{flux-science}:
Einstein--Hilbert, holographic, inflation, and quantum-gravity components; the
\texttt{flux-topology} invariants; the fold), each correspondence was cross-checked
against session measurements (settlement rates, state/verify throughput, the divergence
event record), and every gap found in review was folded back into the text rather than
argued away. The v0.2 corrections — discriminant, gauge language, knot formulation,
composition statement — came from exactly that loop.

\section*{References}

\small
\begin{itemize}[leftmargin=1.2em,itemsep=1pt]
\item \emph{SIGIL — The Variational Chain}, whitepaper v1.1 (2026-07-18) —
\url{https://quillon.xyz/downloads/SIGIL_WHITEPAPER_v1.pdf}
\item \emph{A Closed-Form Delivery Law for Gossip Redundancy in the sigil-top Light
Node} (2026-06-17) — \url{https://quillon.xyz/downloads/sigil-top-delivery-law.pdf}
\item \emph{SIGIL $\times$ QTFT} (2026-05-30) — \texttt{SIGIL\_QTFT\_TOPOLOGY\_v0.md}
\item The SQInstructor: a guide to SQIsign and the Deuring Correspondence —
arXiv:2603.09899
\item High-Throughput GPU Implementation of Dilithium — arXiv:2211.12265
\item A Comparative Analysis of zk-SNARKs and zk-STARKs — arXiv:2512.10020
\end{itemize}

\vspace{0.8em}
\noindent\textit{v0.2 — 2026-07-18. The chain is modeled as a 4-manifold with curvature.
The equation is conjecture; its discipline, and one of its predictions, are not.}

\end{document}
